%\documentclass[PR,cat,default]{P3CTeX}
\input{../../SharedPreamble.tex} %% Relative to the PAC Document loading them all


\def\ISI{Integració de Sistemes d'Informació}
%% Exemple de programació LaTeX avançada
%% Definim un command utilitzant la sintaxi de xparse on:
%% #1 's' Si hi ha un * entre el nom i el primer parametre llavors #1=True
%% #2 'O{PR}'Argument [opcional], si no és present posa PR per defecte.
%% #3 'm' Argument {obligatori}
%% #4 'O{#2}' Argument [opcional] si no és present s'hi posa el mateix valor que #2
%% USAGE: \LoadPreamble*[PR,no-toc]{PR2}{Segona Pràctica}
\NewDocumentCommand{\LoadPreamble}{s O{PR} m O{#2}}{
	% configuració de les metadades utilitzades en les capçaleres peus de pagina i la portada
	\PECTeXconfig{
		data= {
			subj-fullname	= {\ISI},
			subj-shortname	= {ISI},
			subj-code		= {05.597},
			ca-fullname		= {#4},
			ca-shortname	= {#3},
			date			= {\today}
		},
		#2 	%% backdoor to set configs
	}
	% HaCkTeX
	\IfBooleanF{#1}{ %% Si no té * s'executa la portada modificada
		\PECTeXcover{%configuració del disseny de la portada per aquesta asignatura
			outter-backtitle, % vs inner-bactitle
			subtitle,uppertitle, %% both contains same text value from "data={ ca-fullname= this text}"
			subsubtitle, %% Default value "data={ subj-code= this text}"
			text-uppertitle={\vspace*{3ex}\pxGrow{\textcolor{white!60!secondaryColor}{\textbf{ISI [05.597]~}}}}, % Hack default uppertitle text
			text-subsubtitle={2025/2026-2} % Hack default subsubtitle text
		}
	}
	%% superhack: Overwrite no-plagi text:
	\renewcommand{\pxDEFnoplagi}[1]{%%
		\parbox[t][][t]{\linewidth}{
			Aquest es un document d'exemple de com utilitzar \PECTeX reutilitzant el mateix preamble per a totes les PAC i PRs d'una mateixa assignatura.
		}
	}
}
